\documentclass[11pt]{amsart}
\usepackage{amssymb,amsmath,mathtools,booktabs,enumitem}
\usepackage[margin=2.6cm]{geometry}
\usepackage{hyperref}

\newcommand{\stamptime}{\the\numexpr(\time-30)/60\relax:\ifnum\numexpr\time-60*((\time-30)/60)\relax<10 0\fi\the\numexpr\time-60*((\time-30)/60)\relax}
\AddToHook{shipout/foreground}{\put(0.5\paperwidth,-\paperheight+0.55cm){\makebox[0pt]{\tiny\ttfamily unit/doc/\jobname.tex \ --- \ compiled \today, \stamptime}}}

\newtheorem{theorem}{Theorem}[section]
\newtheorem{proposition}[theorem]{Proposition}
\newtheorem{lemma}[theorem]{Lemma}
\newtheorem{corollary}[theorem]{Corollary}
\theoremstyle{definition}
\newtheorem{definition}[theorem]{Definition}
\newtheorem{assumption}{Assumption}
\renewcommand{\theassumption}{A\arabic{assumption}}
\theoremstyle{remark}
\newtheorem{remark}[theorem]{Remark}

\newcommand{\PSh}{\mathrm{PSh}}
\newcommand{\El}{\mathrm{El}}
\newcommand{\Ty}{\mathrm{Ty}}
\newcommand{\Hom}{\mathrm{Hom}}
\newcommand{\Sp}{\mathrm{Sp}}
\newcommand{\Pres}{\mathrm{Pres}}
\newcommand{\Sol}{\mathrm{Sol}}
\newcommand{\Ret}{\mathrm{Ret}}
\newcommand{\yo}{\mathbf{y}}
\newcommand{\pres}[1]{\langle #1\rangle}
\newcommand{\isEquiv}{\mathrm{isEquiv}}
\newcommand{\fib}{\mathrm{fib}}

\title{Duality in the presheaf model:\\ the structure of the proof and its assumptions}
\date{\today}

\begin{document}
\maketitle

\begin{abstract}
We present a proof of the duality axiom of synthetic algebraic geometry in
the presheaf model of Coquand, H\"ofer and Sattler (Theorem~62 of the
version in H\"ofer's licentiate thesis). The proof is decomposed into an
algebraic part in the base, a modal Yoneda lemma, and an identification of
the fibres of the relevant presheaves. Every step after the assumptions is
checked formally. The assumptions are listed with their status: proved for
the concrete model, expected to be provable, or not yet addressed.
\end{abstract}

\section{The statement}

Let $\mathcal C$ be a small category internal to cubical sets, and
$\PSh(\mathcal C)$ the model of presheaves on $\mathcal C$ with values in
cubical sets. Let $D$ be the cobar modality on $\PSh(\mathcal C)$
(H\"ofer's Definition~10), with unit $\tau$; a type $A$ is
\emph{$D$-modal} when $\tau_A\colon A\to DA$ is an equivalence. Let $R$ be
a commutative ring. The objects of $\mathcal C$ are, up to the reversal of
maps, the finitely presented $R$-algebras; $G$ is the generic algebra.

For a presentation $u$ of a $G$-algebra, write $\pres u$ for the presented
$G$-algebra and $\Sp(\pres u) = \Hom_{G}(\pres u, G)$ for its spectrum.

\begin{theorem}[Duality, H\"ofer's Theorem~62]\label{thm:duality}
In the model of $D$-modal types, for every presentation $u$, the transposed
evaluation
\[
  \lambda\,\mathrm{ev}\colon \pres u \longrightarrow G^{\Sp(\pres u)}
\]
is an equivalence.
\end{theorem}

\subsection*{Reading}
All statements are read in the internal language of cubical sets, taken
homotopically: the internal language is univalent type theory, presheaf
fibres are types, and equalities are identifications. For a context
$\Gamma$, an object $x$ of $\mathcal C$ and a point $g\colon \yo x\to\Gamma$,
the \emph{fibre} of a type $A$ over $\Gamma$ at $g$ is
$\fib_A(x,g) = \El(\yo x, A[g])$. A map of types is a \emph{levelwise
equivalence} when it induces an equivalence on all fibres.

Presentations are \emph{syntactic}: a presentation over an algebra $B$ is a
number $n$ of variables and a finite list of polynomial expressions with
coefficients in $B$. (H\"ofer's Construction~55 uses presentations by
elements of free algebras; there is no map from those back to syntax
without choice, and the axioms of synthetic algebraic geometry are stated
for syntactic presentations.)

\section{Assumptions and their status}\label{sec:assumptions}

The proof uses the assumptions below. Their status is one of:
\textbf{proved} (checked for the concrete model), \textbf{expected}
(provable by known means, not yet checked), \textbf{open} (not yet
addressed). The metatheory has univalence and replacement as parameters;
nothing beyond them and the assumptions is used.

\begin{assumption}[the presheaf model]\label{a:cwf}
$\PSh(\mathcal C)$ is a category with families with $\Pi$-types,
identity types and representables $\yo x$, with Yoneda
($\Hom(y,x)\simeq\mathrm{Sub}(\yo y,\yo x)$), the universal property of
context extension, and stability of $\lambda$ under substitution; its
types over a context form a set.
\emph{Status: proved} for the concrete presheaf model on the cubical base,
which also carries $\Sigma$ and identity types with uniform fibrancy.
The equalities of the model are judgemental; reading them as
identifications, and types as a set, is the passage from the external
model to the internal reading, which is \emph{expected}.
\end{assumption}

\begin{assumption}[fibrancy and internal equivalences]\label{a:fib}
Fibrant types are stable under substitution; an internal equivalence acts
as an equivalence on elements; precomposition with an internal
equivalence of contexts is an equivalence on elements of a fibrant type;
a map of types with a two-sided inverse is an internal equivalence.
\emph{Status: expected.} The comparison between internal equivalence and
levelwise equivalence together with naturality in $\mathcal C$ is proved
in both directions for the concrete model; its final packaging is in
progress.
\end{assumption}

\begin{assumption}[the cobar modality]\label{a:D}
$D$ acts on contexts, types and families, with unit $\tau$ natural and
stable under substitution, and with the commutation squares used in the
proof of H\"ofer's Lemma~67.
\emph{Status: proved} for the concrete $D$: $D$ on contexts, types,
substitutions and elements, its functoriality, naturality of $\tau$, and
stability of $D$ under substitution.
\end{assumption}

\begin{assumption}[H\"ofer's Lemma~12 and Proposition~13]\label{a:prop13}
$D$ sends levelwise fibrant types to fibrant types, and a levelwise
equivalence between levelwise fibrant types to an internal equivalence.
\emph{Status: proved} for the concrete $D$ (Lemma~12 in both its
trivial-fibrancy and fibrancy forms, uniformly in $\mathcal C$;
Proposition~13 in its internal form, uniform in $\mathcal C$).
\end{assumption}

\begin{assumption}[$\Pi$-closure]\label{a:modpi}
If $A$ is fibrant and $B$ is $D$-modal over $\Gamma.A$, then $\Pi_A B$
is $D$-modal; and $\Pi$ of fibrant types is fibrant.
\emph{Status:} the fibrancy half is \textbf{proved} for the concrete
model; modality of $\Pi$ is \textbf{expected} (it is part of the inner
model of $D$-modal types, H\"ofer's \S2.2.2).
\end{assumption}

\begin{assumption}[the site]\label{a:site}
The objects of $\mathcal C$ are $R$-algebras with $\mathcal C(y,x)\simeq
\Hom(A_x, A_y)$, functorially; every algebra presented over some $A_x$ has
an object of $\mathcal C$, isomorphic to it under $A_x$.
\emph{Status: open} (H\"ofer's Construction~39 and Lemma~42).
\end{assumption}

\begin{assumption}[syntactic presentations]\label{a:syn}
The points at $x$ of the context of syntactic presentations are the
syntactic presentations over $A_x$, and restriction along a map acts by
change of coefficients. This holds when the natural numbers, lists and
polynomial expressions of the model are computed levelwise.
\emph{Status: open.}
\end{assumption}

\begin{assumption}[presented models]\label{a:pmod}
Restriction to objects $U$ has a right adjoint $R$ on models, natural in
both arguments; free and presented models exist with their universal
properties, in families and stably under substitution; the algebra
structure of a presented model is read fibrewise, and at the identity point
of $\yo x$ its fibre is its algebra of global elements.
\emph{Status:} the transfer of universal properties along $U\dashv R$
(H\"ofer's Construction~51, Lemma~52, Corollary~53) is \textbf{proved};
existence of the models and their fibrewise reading are \textbf{open}.
\end{assumption}

\begin{assumption}[the generic algebra and the spectrum]\label{a:G}
The algebra of global elements of $G$ over $\yo x$ is $A_x$ (H\"ofer's
Remark~46); the spectrum is the type of model maps into $G$, and
evaluation applies such a map.
\emph{Status: open.}
\end{assumption}

\begin{assumption}[modality and fibrancy facts]\label{a:mods}
$\pres u$ and $G$ are $D$-modal; $G$, $\Sp(\pres u)$ and $\pres u$ are
fibrant.
\emph{Status: expected}, from H\"ofer's Theorem~29 ($D$ is a descent data
operation), which is in progress for the concrete model.
\end{assumption}

\section{The algebra in the base}\label{sec:algebra}

This part is ordinary univalent algebra. Fix an $R$-algebra $A$ and a
syntactic presentation $u$ over $A$; let $X_u$ be the presented
$A$-algebra and $p_u\colon A\to X_u$ its structure map. For $h\colon A\to
B$ let $u\cdot h$ be the presentation over $B$ obtained by applying $h$ to
the coefficients.

\begin{lemma}[base change, H\"ofer's Lemma~50]\label{lem:bc}
There is $h^+\colon X_u\to X_{u\cdot h}$ such that
\[
\begin{array}{ccc}
A & \xrightarrow{p_u} & X_u\\
{\scriptstyle h}\downarrow && \downarrow{\scriptstyle h^+}\\
B & \xrightarrow{p_{u\cdot h}} & X_{u\cdot h}
\end{array}
\]
is a pushout of $R$-algebras. The comparison maps between pushouts are
unique, and satisfy the identity and composition laws.
\end{lemma}

\begin{proof}
A cocone $(r,s)$ under $(p_u,h)$ makes $r$ an $A$-algebra map $X_u\to D$
for the structure $s h$; its values on the generators solve $u$, hence
$u\cdot h$ for the structure $s$, which gives a unique $B$-algebra map
$X_{u\cdot h}\to D$. Uniqueness on generators gives the laws.
\end{proof}

At $B=X_u$ and $h=p_u$, the identity of $X_u$ and the identity of $X_u$
form a cocone; the induced map $q_u\colon X_{u\cdot p_u}\to X_u$ is a common
retraction:
\begin{equation}\label{eq:q}
q_u\circ p_u^+ = \mathrm{id}, \qquad q_u\circ p_{u\cdot p_u} = \mathrm{id}.
\end{equation}

\begin{proposition}[base part of H\"ofer's Proposition~58]\label{prop:rep}
For every algebra $B$, the map
\[
\Hom(X_u,B) \longrightarrow \textstyle\sum_{h\colon A\to B}\Ret_B(u\cdot h),
\qquad r\mapsto \bigl(r p_u,\ \text{the retraction induced by }(r,\mathrm{id}_B)\bigr),
\]
is an equivalence, where $\Ret_B(w)$ is the type of retractions of $p_w$.
Moreover the retraction attached to $r$, applied to $h^+(y)$, is $r(y)$.
\end{proposition}

\begin{proof}
For fixed $h$, retractions of $p_{u\cdot h}$ correspond to maps
$r\colon X_u\to B$ with $r p_u=h$, by the pushout property; summing over
$h$, the pairs $(h,r)$ with $r p_u = h$ form $\Hom(X_u,B)$ (singletons).
\end{proof}

\begin{proposition}[base part of H\"ofer's Construction~57]\label{prop:points}
$\Ret_B(w)$ is equivalent to the solutions of $w$ in $B$; under this
equivalence a retraction applied to the value of an expression is the
expression evaluated at the solution. Hence
$\Hom(X_u,B)\simeq\sum_h \Sol(u\cdot h, B)$, and the solution attached to
$r$ is $r$ on the generators of $X_u$.
\end{proposition}

\section{The modal Yoneda lemma}\label{sec:yoneda}

This part uses only Assumptions~\ref{a:cwf}--\ref{a:modpi}.

\begin{lemma}[H\"ofer's Lemma~66]\label{lem:66}
In a commuting diagram of two rows $A_0\to B_0\to C_0$ and $A_1\to B_1\to
C_1$ with vertical maps $a,b,c$, if both horizontal composites and $b$ are
equivalences, so are $a$ and $c$.
\end{lemma}

\begin{proof}
$k = r_0 b^{-1} s_1$ has $k a\simeq r_0 s_0$ and $c k\simeq r_1 s_1$, so
$k$ is an equivalence by $2$-out-of-$6$.
\end{proof}

\begin{lemma}[H\"ofer's Lemma~67]\label{lem:67}
If $f\colon\Theta\to\Gamma$ is a levelwise equivalence of contexts and $Q$
is a fibrant $D$-modal type over $\Gamma$, then restriction along $f$ is an
equivalence $\El(\Gamma,Q)\simeq\El(\Theta,Q[f])$.
\end{lemma}

\begin{proof}
Compare with the same restriction for $DQ$ along $Df$, which is an internal
equivalence by Assumption~\ref{a:prop13}; the horizontal composites are the
action of $\tau_Q$, an equivalence since $Q$ is modal; conclude by
Lemma~\ref{lem:66}.
\end{proof}

\begin{proposition}[H\"ofer's Proposition~69]\label{prop:69}
If $\delta\colon\yo x\to\Delta$ is a levelwise equivalence and $P$ is
fibrant and $D$-modal over $\Delta$, then $u\mapsto u[\delta]$ is an
equivalence $\El(\Delta,P)\simeq\fib_P(x,\delta)$.
\end{proposition}

A \emph{local representation} of a type $A$ over $\Gamma$ assigns to every
point $g$ of $\Gamma$ at $x$ an object $x.g$, a map $p_g\colon x.g\to x$ and
an element $q_g$ of $A$ at $g p_g$, such that $\langle p_g,q_g\rangle\colon
\yo(x.g)\to \yo x.A[g]$ is a levelwise equivalence.

\begin{lemma}[H\"ofer's Lemma~72]\label{lem:72}
For $A$ locally represented and $P$ fibrant and $D$-modal over
$\Gamma.A$, evaluation at $(g p_g,q_g)$ is an equivalence
$\fib_{\Pi_A P}(x,g)\simeq\fib_P(x.g,(gp_g,q_g))$.
\end{lemma}

\begin{proposition}[H\"ofer's Remark~60]\label{prop:60}
For $S$ fibrant over $\Gamma$ and $m\colon S[p]\to P$ over $\Gamma.A$, the
component of $\lambda m\colon S\to\Pi_A P$ at $g$ is an equivalence as soon
as $s\mapsto m_{(gp_g,q_g)}(s\!\restriction_{p_g})$ is one.
\end{proposition}

\begin{proof}
The square relating $\lambda m$, restriction along $p_g$, $m$ at
$(gp_g,q_g)$ and the map of Lemma~\ref{lem:72} commutes (it uses that types
form a set to compare the two points); conclude by $2$-out-of-$3$.
\end{proof}

\begin{proposition}[H\"ofer's Proposition~17]\label{prop:17}
For $D$-modal types, a levelwise equivalence is an internal equivalence,
and conversely.
\end{proposition}

\begin{proof}
The converse is Assumption~\ref{a:fib}. For the direct implication, the
map $e = (p,m)$ of contexts is a levelwise equivalence; Lemma~\ref{lem:67}
gives an inverse $n$ with $n m = \mathrm{id}$ and, by injectivity of
restriction along $e$, $m n=\mathrm{id}$; a two-sided inverse gives an
internal equivalence (Assumption~\ref{a:fib}).
\end{proof}

\section{The fibres of the presentations, the algebras and the spectrum}\label{sec:fibres}

From Assumptions~\ref{a:site}--\ref{a:G}, with $A_x$ the algebra of the
object $x$ and $\Gamma_0$ the context of syntactic presentations:

\begin{proposition}[H\"ofer's Constructions~55--57, syntactic form]\label{prop:fibres}
\begin{enumerate}[label=(\roman*)]
\item The points of $\Gamma_0$ at $x$ are the presentations over $A_x$;
restriction along $f\colon y\to x$ is $u\mapsto u\cdot A_f$.
\item The fibre of $\pres{-}$ at $u$ is $X_u$; restriction along $f$ is
$A_f^+$.
\item The fibre of $G$ at $x$ is $A_x$; restriction is $A_f$.
\item The fibre of $\Sp$ at $u$ is $\Sol(u,A_x)$; restriction along $f$ is
$v\mapsto A_f\circ v$.
\item Evaluation at $(v,b)$ is the universal map $X_u\to A_x$ of the
solution $v$, applied to $b$.
\end{enumerate}
\end{proposition}

\begin{proof}[Proof sketch]
(i) is Assumption~\ref{a:syn}. For (ii), the presented model is computed
fibrewise by the transfer of universal properties along $U\dashv R$
(Construction~51, Lemma~52); at a point it is presented by the restricted
presentation, which over $\yo x$ is a presentation over $A_x$
(Assumption~\ref{a:G}); restriction moves generators to generators, so it
is $A_f^+$ by uniqueness. (iv): a fibre of $\Sp$ is a model map
$\pres u\to G$, which by the universal property is a solution of $u$ in the
global elements of $G$, i.e.\ in $A_x$; restriction is naturality of
evaluation. (v): both sides are algebra maps out of $X_u$ agreeing on the
generators.
\end{proof}

In the last step of H\"ofer's proof of Construction~57 the condition on the
pulled-back family should read $\varepsilon(s_r)\simeq\varepsilon(t_r)$;
the displayed formula there has $s$ on both sides.

\begin{proposition}[H\"ofer's Proposition~58]\label{prop:58}
$\Sp$ is locally represented over $\Gamma_0$: at a point $u$ over $x$, by
the object $x.u$ of $X_u$ (Assumption~\ref{a:site}), $p_u$, and $q_u$.
\end{proposition}

\begin{proof}
At a test object $y$, the points of $\yo(x.u)$ are the algebra maps
$X_u\to A_y$, and the points of $\yo x.\Sp[u]$ are the pairs of a map
$h\colon A_x\to A_y$ and a solution of $u\cdot h$ in $A_y$, by
Proposition~\ref{prop:fibres}. The map $\langle p_u,q_u\rangle$ becomes the
equivalence of Proposition~\ref{prop:points}: the solution attached to $q_u$
restricted along $r$ is $r$ on the generators.
\end{proof}

\section{Proof of duality}\label{sec:proof}

\begin{proof}[Proof of Theorem~\ref{thm:duality}]
$\pres{-}$ is $D$-modal, and so is $G^{\Sp}$ by Assumption~\ref{a:modpi};
by Proposition~\ref{prop:17} it suffices to prove that $\lambda\,\mathrm{ev}$
is a levelwise equivalence. Fix a point $u$ over $x$. By
Proposition~\ref{prop:60} with the local representation of
Proposition~\ref{prop:58}, it suffices that
\[
K\colon\ \fib_{\pres{-}}(x,u)\longrightarrow A_{x.u},\qquad
b\mapsto \mathrm{ev}\bigl(q_u,\ b\!\restriction_{p_u}\bigr)
\]
is an equivalence (through the identification of the fibre of $G$). By
Proposition~\ref{prop:fibres}, $b\!\restriction_{p_u}$ is $p_u^+(b)$ and
$\mathrm{ev}(q_u,-)$ is the universal map of the generators of $X_u$ read in
$A_{x.u}$, which is $q_u$ followed by the isomorphism $X_u\cong A_{x.u}$.
So $K$ is that isomorphism composed with $q_u\circ p_u^+=\mathrm{id}$
\eqref{eq:q}, hence an equivalence.
\end{proof}

The two other presheaf axioms are proved in the same framework: constancy
of $\mathbb N$ (H\"ofer's Theorem~61), with the natural numbers of the model
computed levelwise; projectivity of spectra (H\"ofer's Proposition~76),
which additionally assumes propositional truncation and H\"ofer's
Proposition~34.

\section{What is checked}\label{sec:checked}

The proof above is formalised in the checker Unit
(\texttt{github.com/coquand/unit}), with the assumptions as named
parameters and no axioms beyond the univalence and replacement parameters.

\medskip
\begin{tabular}{@{}p{0.42\textwidth}p{0.52\textwidth}@{}}
\toprule
Content & Development\\
\midrule
\S\ref{sec:algebra} (Lemma~50, $q_u$, Prop.~58 and Constr.~57 in the base)
& \texttt{sag/duality}\\
\S\ref{sec:yoneda} (Lemmas~66, 67, 72, Props.~17, 69, Remark~60)
& \texttt{sag/duality-presheaf}\\
\S\S\ref{sec:fibres}--\ref{sec:proof} (Constrs.~55--57 from the model
assumptions, Prop.~58, Theorem~62, also Theorems~61 and Prop.~76)
& \texttt{sag/duality-levelwise}\\
the concrete $D$ on the cubical base: Definition~10, Lemmas~9, 12, 14, 16
(levelwise), 54, Proposition~13 and Corollary~15 in internal form; the
category with families of uniformly fibrant presheaves on $\mathcal C$
with $\Sigma$, $\Pi$, identity types and function extensionality
& \texttt{sag/cobar}\\
\bottomrule
\end{tabular}

\medskip
\noindent\textbf{Remaining.} (1) Instantiating the abstract assumptions
\ref{a:cwf}--\ref{a:modpi} and \ref{a:mods} by the concrete model, through
the passage from the external model to the internal reading: in progress
(Theorem~29, the comparison of internal and levelwise equivalence).
(2) Assumptions~\ref{a:site}--\ref{a:G}: open; they concern the algebraic
site, the inductive types of the model, and presented models.
(3) Proposition~34, used only for projectivity: to be checked relative to an
interface for propositional truncation in the base.

\end{document}
